% Weekly exercise sheet for Lecture 10; also included in the synthesis sheet.
\ifdefined\ExercisesSynthesis
\else
\documentclass[11pt]{article}
\usepackage[utf8]{inputenc}
\usepackage[T1]{fontenc}
\usepackage[english]{babel}
\usepackage{lmodern,microtype}
\usepackage[margin=0.9in]{geometry}
\usepackage{amsmath,amssymb,amsthm,mathtools}
\usepackage{enumitem,hyperref,needspace}
\newcommand{\E}{\mathbb E}
\newcommand{\Qbb}{\mathbb Q}
\newcommand{\dd}{\,\mathrm d}
\newcommand{\Pb}{\mathbb P}\newcommand{\F}{\mathcal F}\newcommand{\Var}{\operatorname{Var}}
\theoremstyle{definition}
\newtheorem{exercise}{Exercise}[section]
\newlist{exparts}{enumerate}{1}
\setlist[exparts,1]{label=\textbf{(\alph*)},itemsep=0.3em,topsep=0.4em,leftmargin=2em}
\newcommand{\lecturesummary}[1]{\begin{quote}\small\textbf{Focus.} #1\end{quote}}
\title{Stochastic Calculus and Option Pricing\\Lecture 10 Exercises}
\author{}
\date{}
\AddToHook{env/exercise/before}{\Needspace{7\baselineskip}}
\expandafter\newif\csname ifIncludeCore\endcsname\IncludeCoretrue
\expandafter\newif\csname ifIncludeFurther\endcsname\IncludeFurthertrue
\begin{document}
\maketitle
\vspace{-1.2em}
\setcounter{section}{9}
\fi

\section{Week 10 --- American Options and Optimal Stopping}
\lecturesummary{Work with admissible stopping rules, settlement conventions, exercise incentives, and numerical-policy diagnostics.}

\ifIncludeCore

\lecturesummary{Core workload: about 90 minutes, including 10 minutes to check signs, assumptions, and conclusions. Suggested solving times: 10.1: 25 min; 10.2: 15 min; 10.3: 20 min; 10.4: 20 min. These are planning estimates after studying the lecture.}
\setcounter{exercise}{0}
\begin{exercise}[A supermartingale gives a writer's hedge]
Work in a discrete-time, arbitrage-free market with money-market account
\(B_t\). Let \(G_t\) be the amount due if an American claim is exercised at
date \(t\). Suppose a self-financing portfolio has value \(X_t\) and
\(X_t/B_t\) is a nonnegative supermartingale under a risk-neutral measure.
\begin{exparts}
  \item Prove that if \(X_t\geq G_t\) at every exercise date, then the
  initial portfolio value upper-bounds the expected discounted payment for
  every admissible exercise rule.
  \item Explain why checking only the terminal payoff \(G_T\) is
  insufficient for an American claim.
  \item What additional condition would make this superhedge cheapest among
  all portfolios that cover every exercise date? Relate your answer to the
  least supermartingale that dominates the discounted exercise payoff.
\end{exparts}
\end{exercise}

\begin{exercise}[Variation --- Delayed settlement changes the exercise payment]
A holder chooses an exercise time \(\tau\in[t,T]\). The amount
\(g(S_\tau)\geq0\) is fixed at exercise but paid only at
\(\tau+\delta\), with constant delay \(\delta>0\). Rates are constant,
settlement is certain, and exercise opportunities are otherwise unchanged.
\begin{exparts}
\item Express the discounted payment from a chosen \(\tau\) as seen at time \(t\). What is the effective time-\(\tau\) exercise reward?
\item Relate the delayed-settlement American value to the ordinary immediate-settlement value. Does the optimal exercise region change?
\item Now let the known delay depend on the exercise state, \(\delta=\delta(S_\tau)\). Write the new reward and explain why the same constant scaling argument no longer applies.
\item Why would the argument fail if the payment amount were instead recomputed from \(S_{\tau+\delta}\)? State the conditional valuation step then needed at exercise.
\end{exparts}
\end{exercise}

\begin{exercise}[Variation --- A local exercise test with continuous dividends]
Under \(\Qbb\), let
\(\dd S/S=(r-q)\dd t+\sigma\dd W\), with \(r,q\geq0\) and
\(\sigma>0\). Define
\(\mathcal L f=(r-q)s f_s+\tfrac12\sigma^2s^2f_{ss}\).
For an American value, use
\(\max\{g-u,u_t+\mathcal L u-ru\}=0\).
\begin{exparts}
\item Inside a smooth, in-the-money exercise region for a put, substitute \(u=g=K-s\) and derive a necessary inequality involving \(q s\) and \(rK\).
\item Repeat for a call with \(g=s-K\). Explain how dividend yield changes the economic incentive to exercise.
\item Are these local inequalities sufficient to determine the exercise region? Explain what continuation values and boundary conditions still contribute.
\item At a regular put exercise boundary \(s=b(t)\), state value matching and smooth fit. Explain why the boundary must be found as part of the problem.
\end{exparts}
\end{exercise}

\begin{exercise}[Practice --- Audit a Longstaff--Schwartz exercise decision]
At a candidate exercise date, a put has strike 100. A regression fitted
on a separate training sample estimates continuation values
\(\widehat C(80)=18\) and \(\widehat C(90)=12\). On two independent
evaluation paths, the current spots are 80 and 90; their realized
future cash flows, discounted back to the decision date, are 30 and 0.
\begin{exparts}
\item Apply the regression-based rule on the two evaluation paths. State the decisions and the realized cash flows used in the policy evaluation.
\item A programmer instead uses the larger of immediate payoff and the path's realized future cash flow. Compute the resulting two numbers and identify the look-ahead error.
\item Why is a zero realized future cash flow compatible with a positive fitted continuation value? Distinguish a conditional expectation from one observation.
\item Explain why a fixed admissible learned policy has expected value no larger than the optimal stopping value, but neither its finite-sample estimate nor an in-sample optimized estimate is automatically a rigorous lower bound.
\end{exparts}
\end{exercise}
\fi

\ifdefined\ExercisesSynthesis\else
\end{document}
\fi
